\documentclass{report}
\usepackage{amsmath,amssymb}
\setlength{\parindent}{0mm}
\setlength{\parskip}{1em}
\begin{document}
\begin{center}
\rule{6in}{1pt} \
{\large Bryan, D Quaife \\
{\bf Multigrid Applied to Two-Dimensional Vesicle Suspensions}}

ICES \\ University of Texas \\ 201 East 24th St \\ Stop C0200 \\ Austin \\ Texas 78712-1229
\\
{\tt quaife@ices.utexas.edu}\\
George Biros\end{center}

\newcommand{\xx}{{\mathbf{x}}}
\newcommand{\yy}{{\mathbf{y}}}
\newcommand{\ff}{{\mathbf{f}}}
\renewcommand{\SS}{{\mathcal{S}}}
\newcommand{\RR}{{\mathbb{R}}}
\newcommand{\BB}{{\mathcal{B}}}
\newcommand{\DD}{{\mathcal{D}}}
\newcommand{\TT}{{\mathcal{T}}}
\newcommand{\LL}{{\mathcal{L}}}

We are interested in simulating high concentration suspensions of
deformable capsules suspended in a Stokesian fluid. Such capsules are
typically modeled as infinitesimally thin membranes filled with a viscous
fluid. In particular, we are interested in a specific type of capsule
known as a {\em vesicle}. Vesicles are closed inextensible lipid
membranes containing a viscous fluid. The membrane mechanics are
characterized by bending resistance (derived by a Helfrich energy, or a
minimization of the $L^2$ norm of the mean curvature), zero resistance to
shear, and inextensibility (infinite resistance to elongation). Although
we focus on vesicles, the resulting methodologies should be applicable in
a broad range of phenomena in complex fluids and suspensions.

Vesicle evolution dynamics are given by a quasi-static Stokes equation
driven by interface jump conditions derived by a balance of forces on the
membrane. Given the position of the vesicles, the jump conditions can be
evaluated, a Stokes problem with interface conditions can be solved, and
the velocity at the interface advances the vesicle to its new location.
We use a boundary integral formulation that results in a system of
non-linear integro-differential equations for the vesicles's position and
an auxiliary field that is a ``tension'' (it acts as a Lagrange
multiplier to enforce the inextensibility constraint). Below we summarize
the main equations.

Consider a single vesicle $\gamma \in \RR^{2}$ parameterized by the
periodic function $\xx$. The equations governing the motion of the
vesicle are
\begin{align*}
\frac{d\xx}{dt} = \SS[-\xx_{ssss}+(\sigma \xx_{s})_{s}],
\hspace{20pt} \xx_{s} \cdot \frac{d\xx_{s}}{dt} = 0,
\end{align*}
where $\sigma$ is the tension and $\SS[\ff](\xx) =
\int_{\gamma}G(\xx-\yy)\ff(\yy)ds_{\yy}$, where $G(\xx) =
\frac{1}{4\pi\mu}\left(-\log|\xx| + \frac{\xx \otimes
\xx}{|\xx|^{2}}\right)$, is the single-layer potential for Stokes flow.
We introduce the operators $\BB$ (bending), $\TT$ (tension), and $\DD$
(divergence) so that the governing equations are
\begin{align*}
\frac{d\xx}{dt} = \SS[-\BB\xx+\TT\sigma], \hspace{20pt}
\DD\frac{d\xx}{dt} = 0.
\end{align*}

Discretizing in time and linearizing, the vesicle position and tension
are updated by solving
\begin{align}
\left[ \begin{array}{cc}
I + \Delta t \SS\BB & -\Delta t \SS\TT \\
-\DD\SS\BB & \DD\SS\TT
\end{array} \right]
\left[ \begin{array}{c}
\xx^{N+1} \\ \sigma^{N+1}
\end{array} \right] =
\left[ \begin{array}{c}
\xx^{N} \\ 0
\end{array} \right].
\label{e:update}
\end{align}
This system is solved with a matrix-free fast-multipole accelerated GMRES
for the new position and tension. However, the number of iterations
depends on $N$ since the system is ill-conditioned. This motivates the
use of multigrid. In this talk we present the overall scheme and
preliminary results that demonstrate the effectiveness of the proposed
scheme.

As a first step, we are looking for efficient solvers for the Schur
complement of the tension, $\LL:=\DD\SS\TT$. $\LL$ must be inverted for
explicit schemes or it can be used to form block preconditioners for
implicit schemes when solving~\eqref{e:update}. Standard smoothers
applied to integral operators such as $\LL$ generally reduce low
frequencies in the error. We propose using GMRES applied to the high
frequency components as a smoother. In more detail, the smoother for $\LL
\sigma = b$ is $ \sigma^{N+1} = P \sigma^{N} +
\sigma_{H}$, where $\sigma_{H}$ solves $(I-P) \LL \sigma_{H} = (I-P) b -
(I-P)\LL P\sigma^{N}$, and $P$ is the projection onto the low frequency
space.

We use a V-cycle multigrid solver for $\LL$ and report results in
Table~\ref{t:multiSolver}. We see that only a few GMRES iterations are
required to apply the smoother. We set the GMRES tolerance of the
smoother to $10^{-8}$, the coarsest level to $N=8$, and stop the solver
when the relative error is less than $10^{-8}$. We are currently
experimenting with using multigrid to precondition a GMRES solver for
$\LL$.

\begin{table}[h]
\centering
\begin{tabular}{cccc}
N & $\rho$ & $nV$ & GMRES \\
\hline 16 & 0.138 & 10 & 2 \\
32 & 0.365 & 19 & 4 \\
64 & 0.448 & 23 & 8 \\
128 & 0.573 & 35 & 13 \\
\end{tabular}
\caption{\label{t:multiSolver} $N$ is the problem size, $\rho$ is the average
convergence rate, $nV$ is the number of V-cycles, and GMRES is the number
of GMRES iterations required by the smoother at the finest level.}
\end{table}


\end{document}
